% ============================================================================
%  Chapitre 3, partie III : Équations du mouvement
%  Approche Gibbs-Appell
% ============================================================================

% ----------------------------------------------------------------------------
\section{Équations du mouvement - Approche Gibbs-Appell}
% ----------------------------------------------------------------------------

On reprend la sphère de Chaplygin avec la formulation de Gibbs-Appell
du chapitre 2. Au lieu de garder les $5$ coordonnées liées par la
contrainte de roulement, on écrit directement la dynamique avec les
vitesses libres. Il faut noter que cette partie n'est rédigée qu'après la dérivation
du schéma de simulation donc il sert plutôt comme une validation du modèle.

% ----------------------------------------------------------------------------
\subsection{Pseudo-vitesses de la sphère}

\subsubsection{Dimensions}

Après la contrainte holonome $z_G = a$, l'espace de configuration est
de dimension $5$ :
\[
    q = (x, y, \varphi, \theta, \psi)\in\R^5,
    \qquad n = 5
\]
La contrainte de roulement $\mat{A}(q)\,\dot q = 0$ obtenue en II.2.2
compte deux équations, une selon $x$ et une selon $y$. On sépare la
position $r = (x,y)$ et l'orientation $\alpha = (\varphi,\theta,\psi)$ :
\[
    \mat{A}(q) = \begin{pmatrix} \Id_2 & \mat{A}_\alpha(\alpha)
    \end{pmatrix},
    \qquad
    \mat{A}_\alpha(\alpha) = a\begin{pmatrix}
        -s_\psi c_\theta & -c_\psi & 0\\
        c_\psi c_\theta & -s_\psi & 0
    \end{pmatrix}
\]
Le bloc $\Id_2$ montre que $\rang\mat{A}(q) = 2$ pour tout $q$, donc
\[
    k = 2,
    \qquad
    d = n - k = 5 - 2 = 3
\]
Il y a donc $3$ vitesses libres, ce sont celles liées à l'orientation.

\subsubsection{Matrice des vitesses admissibles}

La contrainte s'écrit (on reutilise $r$, et $alpha$ de la partie précédente) :
\[
    \mat{A}(q)\,\dot q = \dot r + \mat{A}_\alpha\,\dot\alpha = 0
    \longleftrightarrow
    \dot r = -\,\mat{A}_\alpha\,\dot\alpha
\]
Une fois $\dot\alpha$ choisi, $\dot r$ est imposé. On prend donc
\[
    \mat{N}(q) = \begin{pmatrix} -\mat{A}_\alpha(\alpha)\\ \Id_3
    \end{pmatrix}\in\mathcal{M}_{5,3}(\R)
\]

\begin{proposition}[Pseudo-vitesses de la sphère]
    Les colonnes de $\mat{N}(q)$ forment une base de
    $\ker\mat{A}(q)$, et les pseudo-vitesses sont les vitesses de
    l'orientation :
    \[
        u = \dot\alpha\in\R^3,
        \qquad
        \dot q = \begin{pmatrix} \dot r\\ \dot\alpha \end{pmatrix}
        = \begin{pmatrix} -\mat{A}_\alpha\,\dot\alpha\\ \dot\alpha
        \end{pmatrix}
    \]
    De même, les déplacements virtuels admissibles s'écrivent
    \[
        \delta q = \mat{N}(q)\,\xi
        \longleftrightarrow
        \delta r = -\,\mat{A}_\alpha\,\xi,
        \qquad
        \delta\alpha = \xi,
        \quad \text{avec} \quad \xi\in\R^3
    \]
\end{proposition}

\begin{remarque}[Vitesses indépendantes]
    Les trois composantes de $u = (\dot\varphi,\dot\theta,\dot\psi)$
    sont indépendantes : on peut les choisir librement, et toute
    valeur de $u$ donne une vitesse admissible. Ce n'est pas le cas de
    $\dot x$ et $\dot y$, qui ne sont pas libres puisqu'elles se
    déduisent de $u$ par le roulement :
    \[
        \begin{pmatrix} \dot x\\ \dot y \end{pmatrix}
        = -\,\mat{A}_\alpha(\alpha)\,u
    \]
    On y retrouve donc les mêmes configurations que l'approche d'EL.
\end{remarque}

% ----------------------------------------------------------------------------
\subsection{Fonction de Gibbs-Appell de la sphère}

On applique la fonction de Gibbs-Appell à la sphère :
\[
    \mathcal{G} = \frac{1}{2}\,m\,\|\vect{a}_G\|^2
    + \frac{1}{2}\,\dot{\vect{\omega}}^T\mat{I}_b\,\dot{\vect{\omega}}
    + \dot{\vect{\omega}}\cdot\bigl(\vect{\omega}\wedge
    \mat{I}_b\vect{\omega}\bigr)
\]
à des termes indépendants de $\dot u$ près. Il reste à exprimer
$\vect{a}_G$ et $\dot{\vect{\omega}}$ en fonction de
$(\alpha,\dot\alpha,\ddot\alpha)$.

\subsubsection{Accélération du centre}

Comme $z_G = a$, $\vect{a}_G = (\ddot x,\ddot y,0)$. En dérivant la
contrainte $\dot r = -\,\mat{A}_\alpha\,\dot\alpha$,
\[
    \ddot r = -\,\bigl(\mat{A}_\alpha\,\ddot\alpha
    + \dot{\mat{A}}_\alpha\,\dot\alpha\bigr)
    \longrightarrow
    \|\vect{a}_G\|^2 = \bigl\|\mat{A}_\alpha\,\ddot\alpha
    + \dot{\mat{A}}_\alpha\,\dot\alpha\bigr\|^2
\]

\subsubsection{Accélération angulaire}

On note $\mat{E}(\alpha)$ la matrice de I.4.3 telle que
\[
    \vect{\omega} = \mat{E}(\alpha)\,\dot\alpha,
    \qquad
    \mat{E}(\alpha) = \begin{pmatrix}
        1 & 0 & -s_\theta\\
        0 & c_\varphi & s_\varphi c_\theta\\
        0 & -s_\varphi & c_\varphi c_\theta
    \end{pmatrix}
\]
En dérivant,
\[
    \dot{\vect{\omega}} = \mat{E}\,\ddot\alpha + \dot{\mat{E}}\,\dot\alpha
\]
on  a :
\begin{proposition}[Fonction de Gibbs-Appell de la sphère]
    À des termes indépendants de $\ddot\alpha$ près,
    \[
        \mathcal{G}(\alpha,\dot\alpha,\ddot\alpha)
        = \frac{1}{2}\,m\,\bigl\|\mat{A}_\alpha\,\ddot\alpha
        + \dot{\mat{A}}_\alpha\,\dot\alpha\bigr\|^2
        + \frac{1}{2}\,\bigl(\mat{E}\,\ddot\alpha
        + \dot{\mat{E}}\,\dot\alpha\bigr)^T\mat{I}_b\,
        \bigl(\mat{E}\,\ddot\alpha + \dot{\mat{E}}\,\dot\alpha\bigr)
        + \bigl(\mat{E}\,\ddot\alpha + \dot{\mat{E}}\,\dot\alpha\bigr)
        \cdot\bigl(\vect{\omega}\wedge\mat{I}_b\vect{\omega}\bigr)
    \]
    avec $\vect{\omega} = \mat{E}\,\dot\alpha$.
\end{proposition}

% ----------------------------------------------------------------------------
\subsection{Équations de Gibbs-Appell}

Le potentiel de pesanteur $V = m g a$ est constant et aucune force non
conservative n'est appliquée, donc
\[
    \mathcal{Q} = \mat{N}(q)^T\left(-\pd{V}{q}\right) = 0
\]
Les équations de Gibbs-Appell s'écrivent donc
\[
    \pd{\mathcal{G}}{\ddot\alpha_j} = 0,
    \qquad \alpha_j\in\{\varphi,\theta,\psi\}
\]
En dérivant $\mathcal{G}$, on a : 
\[
    \pd{\mathcal{G}}{\ddot\alpha_j}
    = m\,\pd{\ddot r}{\ddot\alpha_j}\cdot\ddot r
    + \pd{\dot{\vect{\omega}}}{\ddot\alpha_j}\cdot
    \bigl(\mat{I}_b\,\dot{\vect{\omega}}
    + \vect{\omega}\wedge\mat{I}_b\vect{\omega}\bigr)
\]
avec, d'après $\ddot r = -\,\bigl(\mat{A}_\alpha\,\ddot\alpha,
    + \dot{\mat{A}}_\alpha\,\dot\alpha\bigr)$ et $\dot{\vect{\omega}} = \mat{E}\,\ddot\alpha + \dot{\mat{E}}\,\dot\alpha$, on a : 

\[
    \pd{\ddot r}{\ddot\alpha_j} = -\,(\text{$j$-ième colonne de }
    \mat{A}_\alpha),
    \qquad
    \pd{\dot{\vect{\omega}}}{\ddot\alpha_j} = \text{$j$-ième colonne
    de } \mat{E}
\]
On note $\vect{\Gamma} = \mat{I}_b\,\dot{\vect{\omega}}
+ \vect{\omega}\wedge\mat{I}_b\vect{\omega}$. Avec
$\mat{I}_b = \mathrm{diag}(I_1,I_2,I_3)$,
\[
    \vect{\Gamma} = \begin{pmatrix}
        I_1\,\dot\omega_1 + (I_3 - I_2)\,\omega_2\,\omega_3\\
        I_2\,\dot\omega_2 + (I_1 - I_3)\,\omega_3\,\omega_1\\
        I_3\,\dot\omega_3 + (I_2 - I_1)\,\omega_1\,\omega_2
    \end{pmatrix}
\]

\subsubsection{Roulis : \texorpdfstring{$\varphi$}{phi}}

La première colonne de $\mat{E}$ est $(1,0,0)$ et celle de
$\mat{A}_\alpha$ est $a\,(-s_\psi c_\theta,\, c_\psi c_\theta)$, donc
\[
    \pd{\dot{\vect{\omega}}}{\ddot\varphi} = \begin{pmatrix}
        1\\ 0\\ 0
    \end{pmatrix},
    \qquad
    \pd{\ddot r}{\ddot\varphi} = a\,c_\theta\begin{pmatrix}
        s_\psi\\ -c_\psi
    \end{pmatrix}
\]
L'équation en $\varphi$ est donc
\[
    I_1\,\dot\omega_1 + (I_3 - I_2)\,\omega_2\,\omega_3
    + m\,a\,c_\theta\,(s_\psi\,\ddot x - c_\psi\,\ddot y) = 0
\]

\subsubsection{Tangage : \texorpdfstring{$\theta$}{theta}}

La deuxième colonne de $\mat{E}$ est $(0, c_\varphi, -s_\varphi)$ et
celle de $\mat{A}_\alpha$ est $a\,(-c_\psi, -s_\psi)$, donc
\[
    \pd{\dot{\vect{\omega}}}{\ddot\theta} = \begin{pmatrix}
        0\\ c_\varphi\\ -s_\varphi
    \end{pmatrix},
    \qquad
    \pd{\ddot r}{\ddot\theta} = a\begin{pmatrix}
        c_\psi\\ s_\psi
    \end{pmatrix}
\]
L'équation en $\theta$ est donc
\[
    c_\varphi\bigl(I_2\,\dot\omega_2 + (I_1 - I_3)\,\omega_3\,\omega_1
    \bigr)
    - s_\varphi\bigl(I_3\,\dot\omega_3 + (I_2 - I_1)\,\omega_1\,
    \omega_2\bigr)
    + m\,a\,(c_\psi\,\ddot x + s_\psi\,\ddot y) = 0
\]

\subsubsection{Azimut : \texorpdfstring{$\psi$}{psi}}

La troisième colonne de $\mat{E}$ est $(-s_\theta, s_\varphi c_\theta,
c_\varphi c_\theta)$ et celle de $\mat{A}_\alpha$ est nulle, donc
\[
    \pd{\dot{\vect{\omega}}}{\ddot\psi} = \begin{pmatrix}
        -s_\theta\\ s_\varphi c_\theta\\ c_\varphi c_\theta
    \end{pmatrix},
    \qquad
    \pd{\ddot r}{\ddot\psi} = 0
\]
L'équation en $\psi$ est donc
\[
    -\,s_\theta\bigl(I_1\,\dot\omega_1 + (I_3 - I_2)\,\omega_2\,
    \omega_3\bigr)
    + s_\varphi c_\theta\bigl(I_2\,\dot\omega_2 + (I_1 - I_3)\,
    \omega_3\,\omega_1\bigr)
    + c_\varphi c_\theta\bigl(I_3\,\dot\omega_3 + (I_2 - I_1)\,
    \omega_1\,\omega_2\bigr) = 0
\]

% ----------------------------------------------------------------------------
\subsection{Résumé}

Les inconnues sont $\alpha = (\varphi,\theta,\psi)$ seulement, soit
$3$ fonctions du temps pour $3$ équations.

\begin{theoreme}[Équations de Gibbs-Appell de la sphère de Chaplygin]
    Avec $\vect{\omega} = \mat{E}(\alpha)\,\dot\alpha$ et
    \[
        \ddot r = -\,\bigl(\mat{A}_\alpha\,\ddot\alpha
        + \dot{\mat{A}}_\alpha\,\dot\alpha\bigr)
    \]
    le mouvement de la sphère de Chaplygin vérifie
    \begin{align*}
        &I_1\,\dot\omega_1 + (I_3 - I_2)\,\omega_2\,\omega_3
            + m\,a\,c_\theta\,(s_\psi\,\ddot x - c_\psi\,\ddot y) = 0\\
        &c_\varphi\bigl(I_2\,\dot\omega_2 + (I_1 - I_3)\,\omega_3\,
            \omega_1\bigr)
            - s_\varphi\bigl(I_3\,\dot\omega_3 + (I_2 - I_1)\,
            \omega_1\,\omega_2\bigr)
            + m\,a\,(c_\psi\,\ddot x + s_\psi\,\ddot y) = 0\\
        &-\,s_\theta\bigl(I_1\,\dot\omega_1 + (I_3 - I_2)\,\omega_2\,
            \omega_3\bigr)
            + s_\varphi c_\theta\bigl(I_2\,\dot\omega_2
            + (I_1 - I_3)\,\omega_3\,\omega_1\bigr)
            + c_\varphi c_\theta\bigl(I_3\,\dot\omega_3
            + (I_2 - I_1)\,\omega_1\,\omega_2\bigr) = 0
    \end{align*}
\end{theoreme}